\documentclass[twoside]{article}

\usepackage[preprint]{aistats2027}
\usepackage[round,authoryear]{natbib}
\usepackage{amsmath,amssymb,amsthm}
\usepackage{mathtools}
\usepackage{microtype}
\usepackage{url}

\renewcommand{\bibname}{References}
\renewcommand{\bibsection}{\subsubsection*{\bibname}}
\bibliographystyle{apalike}

\newtheorem{theorem}{Theorem}
\newtheorem{lemma}{Lemma}
\newtheorem{corollary}{Corollary}
\newcommand{\E}{\mathbb E}
\newcommand{\calone}{\operatorname{Cal}_1}
\newcommand{\caltwo}{\operatorname{Cal}_2}
\newcommand{\Reg}{\operatorname{Reg}}
\newcommand{\simplex}{\Delta_K}
\newcommand{\1}{\mathbf 1}

\begin{document}
\runningauthor{}
\onecolumn

\aistatstitle{Supplement: The Price of Fidelity in Multiclass Online Recalibration}

\section{Explicit Simplex Net}

Fix $h\in(0,1)$ and put $m=\lceil\sqrt K/h\rceil$. Define
\[
 N_h=\left\{z/m:z\in\mathbb Z_+^K,\ \sum_{i=1}^Kz_i=m\right\}.
\]
For $p\in\simplex$, first set $z_i=\lfloor mp_i\rfloor$. The deficit
$m-\sum_i z_i$ is an integer between zero and $K-1$. Add one to that many
coordinates with largest fractional parts. The resulting vector still sums to
$m$, and every coordinate differs from $mp_i$ by less than one. Thus the map
$R_h(p)=z/m$ satisfies
\[
 \|R_h(p)-p\|_2<\frac{\sqrt K}{m}\leq h.
\]
Stars and bars gives
\[
 |N_h|={m+K-1\choose K-1}\leq C_Km^{K-1}\leq C_Kh^{-(K-1)}.
\]
This proves all net properties used in Lemma 1 of the main paper.

\section{Strongly Proper Losses}

A multiclass loss $\ell$ is $\lambda$-strongly proper in Euclidean norm if
for every $p,q\in\simplex$,
\[
 \E_{Y\sim q}[\ell(p,Y)-\ell(q,Y)]
 \geq\lambda\|p-q\|_2^2.
\]
The Brier loss in the main paper has equality with $\lambda=1$.

\begin{theorem}[Strongly proper extension]\label{thm:strong}
Fix $K\geq2$ and $\lambda>0$. There are constants
$c_{K,\lambda},C_{K,\lambda},\beta_{K,\lambda}>0$ such that, on the truthful
product instance from the main paper, every learner with
\[
 \E\left[\frac1T\sum_{t=1}^T
 \bigl(\ell(p_t,y_t)-\ell(q_t,y_t)\bigr)\right]\leq\beta
\]
for a $\lambda$-strongly proper loss obeys
\[
 \E\calone(p,y)\geq
 c_{K,\lambda}\min\left\{1,
 (T\beta^{(K-1)/2})^{-1/2}\right\}
 -C_{K,\lambda}\sqrt\beta.
\]
\end{theorem}

\begin{proof}
Conditional strong propriety and truthfulness give
\[
 \E\sum_{t=1}^T\|p_t-q_t\|_2^2\leq\frac{\beta T}{\lambda}.
\]
Choose $h=c\sqrt{\beta/\lambda}$. The rounding argument then gives
$\E B\leq C_Kh^2T$. Every remaining lemma depends only on this movement
bound, the source lattice, and the interior label variance. Therefore the proof
in the main paper gives
\[
 \E\calone(p,y)\geq
 c_K\min\{1,(Th^{K-1})^{-1/2}\}-C_Kh.
\]
Substitution absorbs powers of $\lambda$ into the displayed constants.
\end{proof}

The theorem requires only a lower curvature bound. It does not require
smoothness, mixability, or bounded loss values beyond integrability of the
stated expected excess loss.

\section{Intrinsic Hint Dimension}

We prove Proposition 1 of the main paper. Fix $1\leq r\leq K-1$. Choose an
$r$-dimensional affine box $Q_r$ inside the simplex by varying the first $r$
coordinates in a sufficiently small positive interval, holding coordinates
$r+1$ through $K-1$ at positive constants, and using the last coordinate as
the remainder. A spacing-$h$ grid in $Q_r$ has
$M=\Theta_{K,r}(h^{-r})$ points and the same interior variance bound used in
the main proof.

Let $P_{Q_r}$ be Euclidean projection onto the closed convex box $Q_r$, which
has an $h$-net $N_{h,r}$ of size at most $C_{K,r}h^{-r}$. Define $r_t$ by
projecting $p_t$ and rounding the projection to this net. Since every $q_t$
lies in $Q_r$, the projection inequality gives
\[
 \|P_{Q_r}(p_t)-q_t\|_2\leq\|p_t-q_t\|_2
\]
and
\[
 \operatorname{dist}(p_t,Q_r)\leq\|p_t-q_t\|_2.
\]
The grouping argument from the main rounding lemma now gives
\[
 \calone(r,y)\leq\calone(p,y)
 +\frac1T\sum_{t=1}^T\operatorname{dist}(p_t,Q_r)+h.
\]
Under expected fidelity at most $\beta$, Jensen and Cauchy--Schwarz bound the
expected projection term by $\sqrt\beta$. The rounded movement still obeys
\[
 \E\sum_{t=1}^T\|r_t-q_t\|_2^2\leq C_{K,r}(\beta+h^2)T.
\]

All rounded bins now lie in the $r$-dimensional affine space. Repeat the
martingale proof with
\[
 s=1+\frac2r,
 \qquad
 a=\frac{2s-1}{2s-2},
 \qquad
 b=\frac1{2s-2}.
\]
The source ball count and occupancy transport argument use dimension $r$.
The net cardinality is $O_{K,r}(h^{-r})$. Thus the global occupancy moment is
$O_{K,r}(T^sh^2)$ in the long-horizon regime, and the same exponent
cancellation gives
\[
 \E\calone(r,y)
 \geq c_{K,r}\min\{1,(Th^r)^{-1/2}\}-C_{K,r}h.
\]
Returning from $r_t$ to $p_t$ costs at most $\sqrt\beta+h$ in expectation.
Setting $h=c\sqrt\beta$ proves Proposition 1.

\section{Fractional Martingale Moment}

We expand the interpolation step in Lemma 3 of the main paper. Fix an output
bin $v$ and let $Z=\mu_v$. Put
\[
 s=1+\frac2d,
 \qquad
 \theta=\frac{s-1}{2s-1}.
\]
The interpolation identity
\[
 \frac12=\frac{\theta}{1}+\frac{1-\theta}{2s}
\]
and log-convexity of $L_p$ norms give
\[
 \|Z\|_2\leq\|Z\|_1^\theta\|Z\|_{2s}^{1-\theta}.
\]
Rearranging and raising to the required powers yields
\[
 \E|Z|
 \geq
 \frac{(\E Z^2)^a}{(\E|Z|^{2s})^b},
 \qquad
 a=\frac{2s-1}{2s-2},
 \qquad
 b=\frac1{2s-2}.
\]
The interior variance and Burkholder--Rosenthal bounds therefore imply
\[
 \E|\mu_v|
 \geq c_K\frac{(\E\nu_v)^a}{(\E\nu_v^s)^b}.
\]

For completeness, the Burkholder--Rosenthal inequality used here states that
for a discrete martingale $M_t$ with differences $D_t$ and every real
$p\geq2$,
\[
 \E\max_{t\leq T}|M_t|^p
 \leq C_p\left(
 \E\langle M\rangle_T^{p/2}
 +\E\max_{t\leq T}|D_t|^p\right).
\]
This is the form quoted by \citet[Lemma 1]{collina2026optimal}, who attribute
it to \citet{burkholder1973distribution}. In our application,
$D_t=\1\{r_t=v\}(q_{t,1}-y_{t,1})$. Predictability gives
\[
 \langle M\rangle_T
 =\sum_t\1\{r_t=v\}q_{t,1}(1-q_{t,1})\leq\nu_v.
\]
Also,
$\max_t|D_t|^{2s}\leq\1\{\nu_v>0\}\leq\nu_v^s$. Hence
$\E|\mu_v|^{2s}\leq C_K\E\nu_v^s$ for the noninteger value
$2s=2+4/d$.

To combine bins, define
\[
 A_v=\E\nu_v,
 \qquad B_v=\E\nu_v^s,
 \qquad X_v=A_v^a/B_v^b.
\]
Since $a-1=b$, H\"older with conjugate exponents $a$ and $a/(a-1)$ gives
\[
 \sum_v A_v
 =\sum_vX_v^{1/a}B_v^{b/a}
 \leq\left(\sum_vX_v\right)^{1/a}
 \left(\sum_vB_v\right)^{b/a}.
\]
As $\sum_vA_v=T$, rearrangement gives
\[
 \sum_v\E|\mu_v|
 \geq c_K\frac{T^a}{(\sum_v\E\nu_v^s)^b}.
\]

\section{Detailed Lattice Transport Proof}

Let the source chart be
$q(x)=(x_1,\ldots,x_d,1-\sum_jx_j)$. For any $x,x'\in\mathbb R^d$,
\[
 \|q(x)-q(x')\|_2^2
 =\|x-x'\|_2^2+\left(\sum_j(x_j-x'_j)\right)^2
 \geq\|x-x'\|_2^2.
\]
Therefore a Euclidean ball of radius $R$ in the simplex intersects no more
source points than a radius-$R$ ball intersects in the spacing-$h$ Cartesian
lattice. A coordinate cube cover gives at most
$C_K(1+R/h)^d$ distinct points. Since each source point occurs at most $L$
times, the occurrence count is at most
\begin{equation}\label{eq:ballcount}
 C_KL(1+R/h)^d.
\end{equation}

Fix an output $v$ and sort the distances of assigned source occurrences as
$D_1\leq\cdots\leq D_{\nu_v}$. The center $v$ need not lie on the lattice and
may lie on the simplex boundary. Equation~\eqref{eq:ballcount} still applies.
There are constants $c_K,C_K$ such that for every $i>C_KL$,
\[
 D_i\geq c_Kh(i/L)^{1/d}.
\]
Consequently,
\begin{align*}
 b_v=\sum_{i=1}^{\nu_v}D_i^2
 &\geq c_Kh^2L^{-2/d}
 \sum_{i=C_KL+1}^{\nu_v}i^{2/d}\\
 &\geq c_Kh^2L^{-2/d}
 \left(\nu_v^{1+2/d}-C_KL^{1+2/d}\right).
\end{align*}
Moving the baseline term to the other side proves the occupancy transport
inequality. The proof uses only an upper bound on source multiplicity, so it
also covers a horizon that stops partway through the final lattice cycle.

Summing the transport inequality gives
\[
 \sum_v\nu_v^s
 \leq C_K\left(
 N_{\mathrm{occ}}L^s+L^{2/d}B/h^2\right),
\]
where $N_{\mathrm{occ}}$ is the number of occupied output bins. If $T<M$, then
$L=1$ and $N_{\mathrm{occ}}\leq T$. If $T\geq M$, then
$L\leq C_KTh^d$ and $N_{\mathrm{occ}}\leq|N_h|\leq C_Kh^{-d}$. These two cases
give the global occupancy lemma after taking expectations.

\section{Expected Upper-Bound Conversion}

We detail Proposition 2 of the main paper. Theorem 17 of
\citet{chen2026calibeating} gives a pure-forecast algorithm with expected
cumulative Brier regret $O(\eta^2T)$ relative to any reference algorithm. It
also gives, with probability at least $1-\delta$,
\[
 K_T\leq\widetilde O_K\left(
 \sqrt T+\eta^{-d}+\eta^2T+\log(1/\delta)
 \right),
\]
where $K_T=T\caltwo^2$ and the notation hides fixed-$K$ and logarithmic
factors. Choose the reference algorithm that outputs the revealed hint $q_t$.
Choosing $\delta=T^{-2}$ and using the deterministic bound $K_T\leq2T$ on the
failure event gives
\[
 \E K_T\leq\widetilde O_K(\sqrt T+\eta^{-d}+\eta^2T).
\]
Jensen and $\sqrt{x+y+z}\leq\sqrt x+\sqrt y+\sqrt z$ yield
\[
 \E\caltwo
 \leq\sqrt{\E K_T/T}
 \leq\widetilde O_K\left(
 T^{-1/4}+\frac1{\sqrt{T\eta^d}}+\eta\right).
\]
For $d\geq2$, set $\eta=\varepsilon$ and
$T=\widetilde\Omega_K(\varepsilon^{-(d+2)})$. Then
\[
 T^{-1/4}\leq\widetilde O_K(\varepsilon^{(d+2)/4})
 \leq\widetilde O_K(\varepsilon),
\]
and the other two terms are also $\widetilde O_K(\varepsilon)$. This explains
why the multiclass upper theorem covers $K\geq3$. The binary case uses the
sharper direct construction of \citet{hu2026optimal}.

\section{Deterministic Adversarial Realization}

The lower construction randomizes labels, while the online guarantee in the
main paper quantifies over deterministic sequences and averages only over the
learner seed. Suppose an algorithm met both target guarantees for every
deterministic label realization under the fixed hint sequence. Averaging each
guarantee over the independent categorical labels would preserve both
inequalities. The price-of-fidelity theorem rules this out below the stated
horizon. Therefore at least one deterministic label realization violates at
least one target guarantee. This is the standard averaging form of the
probabilistic method and requires no adaptive adversary.

\section{Calibration Norms}

For weights $w_v=n_v/T$ that sum to one, Cauchy--Schwarz gives
\[
 \calone
 =\sum_vw_v\|v-\rho_v\|_2
 \leq\left(\sum_vw_v\|v-\rho_v\|_2^2\right)^{1/2}
 =\caltwo.
\]
Thus every lower bound on expected $\calone$ applies directly to expected
$\caltwo$. In the binary case, identifying a scalar probability $p$ with
$(p,1-p)$ changes both norms by only a factor $\sqrt2$. No rate exponent is
affected.

\bibliography{references}

\end{document}
